% Options for packages loaded elsewhere
% Options for packages loaded elsewhere
\PassOptionsToPackage{unicode}{hyperref}
\PassOptionsToPackage{hyphens}{url}
%
\documentclass[
  11pt,
  letterpaper,
  DIV=11,
  numbers=noendperiod,
  twoside]{scrartcl}
\usepackage{xcolor}
\usepackage[left=1.1in, right=1in, top=0.8in, bottom=0.8in,
paperheight=9.5in, paperwidth=7in, includemp=TRUE, marginparwidth=0in,
marginparsep=0in]{geometry}
\usepackage{amsmath,amssymb}
\setcounter{secnumdepth}{3}
\usepackage{iftex}
\ifPDFTeX
  \usepackage[T1]{fontenc}
  \usepackage[utf8]{inputenc}
  \usepackage{textcomp} % provide euro and other symbols
\else % if luatex or xetex
  \usepackage{unicode-math} % this also loads fontspec
  \defaultfontfeatures{Scale=MatchLowercase}
  \defaultfontfeatures[\rmfamily]{Ligatures=TeX,Scale=1}
\fi
\usepackage{lmodern}
\ifPDFTeX\else
  % xetex/luatex font selection
  \setmathfont[]{Garamond-Math}
\fi
% Use upquote if available, for straight quotes in verbatim environments
\IfFileExists{upquote.sty}{\usepackage{upquote}}{}
\IfFileExists{microtype.sty}{% use microtype if available
  \usepackage[]{microtype}
  \UseMicrotypeSet[protrusion]{basicmath} % disable protrusion for tt fonts
}{}
\usepackage{setspace}
% Make \paragraph and \subparagraph free-standing
\makeatletter
\ifx\paragraph\undefined\else
  \let\oldparagraph\paragraph
  \renewcommand{\paragraph}{
    \@ifstar
      \xxxParagraphStar
      \xxxParagraphNoStar
  }
  \newcommand{\xxxParagraphStar}[1]{\oldparagraph*{#1}\mbox{}}
  \newcommand{\xxxParagraphNoStar}[1]{\oldparagraph{#1}\mbox{}}
\fi
\ifx\subparagraph\undefined\else
  \let\oldsubparagraph\subparagraph
  \renewcommand{\subparagraph}{
    \@ifstar
      \xxxSubParagraphStar
      \xxxSubParagraphNoStar
  }
  \newcommand{\xxxSubParagraphStar}[1]{\oldsubparagraph*{#1}\mbox{}}
  \newcommand{\xxxSubParagraphNoStar}[1]{\oldsubparagraph{#1}\mbox{}}
\fi
\makeatother


\usepackage{longtable,booktabs,array}
\newcounter{none} % for unnumbered tables
\usepackage{calc} % for calculating minipage widths
% Correct order of tables after \paragraph or \subparagraph
\usepackage{etoolbox}
\makeatletter
\patchcmd\longtable{\par}{\if@noskipsec\mbox{}\fi\par}{}{}
\makeatother
% Allow footnotes in longtable head/foot
\IfFileExists{footnotehyper.sty}{\usepackage{footnotehyper}}{\usepackage{footnote}}
\makesavenoteenv{longtable}
\usepackage{graphicx}
\makeatletter
\newsavebox\pandoc@box
\newcommand*\pandocbounded[1]{% scales image to fit in text height/width
  \sbox\pandoc@box{#1}%
  \Gscale@div\@tempa{\textheight}{\dimexpr\ht\pandoc@box+\dp\pandoc@box\relax}%
  \Gscale@div\@tempb{\linewidth}{\wd\pandoc@box}%
  \ifdim\@tempb\p@<\@tempa\p@\let\@tempa\@tempb\fi% select the smaller of both
  \ifdim\@tempa\p@<\p@\scalebox{\@tempa}{\usebox\pandoc@box}%
  \else\usebox{\pandoc@box}%
  \fi%
}
% Set default figure placement to htbp
\def\fps@figure{htbp}
\makeatother





\setlength{\emergencystretch}{3em} % prevent overfull lines

\providecommand{\tightlist}{%
  \setlength{\itemsep}{0pt}\setlength{\parskip}{0pt}}



 


\setlength\heavyrulewidth{0ex}
\setlength\lightrulewidth{0ex}
\usepackage[automark]{scrlayer-scrpage}
\clearpairofpagestyles
\cehead{
  Brian Weatherson
  }
\cohead{
  Comments on Orthogonal Knowledge
  }
\ohead{\bfseries \pagemark}
\cfoot{}
\makeatletter
\newcommand*\NoIndentAfterEnv[1]{%
  \AfterEndEnvironment{#1}{\par\@afterindentfalse\@afterheading}}
\makeatother
\NoIndentAfterEnv{itemize}
\NoIndentAfterEnv{enumerate}
\NoIndentAfterEnv{description}
\NoIndentAfterEnv{quote}
\NoIndentAfterEnv{equation}
\NoIndentAfterEnv{longtable}
\NoIndentAfterEnv{abstract}
\renewenvironment{abstract}
 {\vspace{-1.25cm}
 \quotation\small\noindent\emph{Abstract}:}
 {\endquotation}
\newfontfamily\tfont{EB Garamond}
\addtokomafont{disposition}{\rmfamily}
\addtokomafont{descriptionlabel}{\rmfamily}
\addtokomafont{title}{\normalfont\itshape}
\let\footnoterule\relax

\makeatletter
\renewcommand{\@maketitle}{%
  \newpage
  \null
  \vskip 2em%
  \begin{center}%
  \let \footnote \thanks
    {\itshape\huge\@title \par}%
    \vskip 0.5em%  % Reduced from default
    {\large
      \lineskip 0.3em%  % Reduced from default 0.5em
      \begin{tabular}[t]{c}%
        \@author
      \end{tabular}\par}%
    \vskip 0.5em%  % Reduced from default
    {\large \@date}%
  \end{center}%
  \par
  }
\makeatother
\RequirePackage{lettrine}

\renewenvironment{abstract}
 {\quotation\small\noindent\emph{Abstract}:}
 {\endquotation\vspace{-0.02cm}}

\setmainfont{EB Garamond Math}[
  BoldFont = {EB Garamond SemiBold},
  ItalicFont = {EB Garamond Italic},
  RawFeature = {+smcp},
]

\newfontfamily\scfont{EB Garamond Regular}[RawFeature=+smcp]
\renewcommand{\textsc}[1]{{\scfont #1}}

\renewcommand{\LettrineTextFont}{\scfont}
% Number theorem-like environments straight through (Theorem 1, Theorem 2, ...)
% rather than by section (1.1, 1.2, ...), which is Quarto's LaTeX default.
% Guarded, since Quarto only defines a counter for environments the document uses.
\makeatletter
\AtBeginDocument{%
  \@for\bw@thm:={theorem,lemma,corollary,proposition,conjecture,definition,%
                 example,exercise,refremark,refsolution}\do{%
    \@ifundefined{c@\bw@thm}{}{\expandafter\counterwithout\expandafter{\bw@thm}{section}}%
  }%
}
\makeatother
\KOMAoption{captions}{tableheading}
\makeatletter
\@ifpackageloaded{caption}{}{\usepackage{caption}}
\AtBeginDocument{%
\ifdefined\contentsname
  \renewcommand*\contentsname{Table of contents}
\else
  \newcommand\contentsname{Table of contents}
\fi
\ifdefined\listfigurename
  \renewcommand*\listfigurename{List of Figures}
\else
  \newcommand\listfigurename{List of Figures}
\fi
\ifdefined\listtablename
  \renewcommand*\listtablename{List of Tables}
\else
  \newcommand\listtablename{List of Tables}
\fi
\ifdefined\figurename
  \renewcommand*\figurename{Figure}
\else
  \newcommand\figurename{Figure}
\fi
\ifdefined\tablename
  \renewcommand*\tablename{Table}
\else
  \newcommand\tablename{Table}
\fi
}
\@ifpackageloaded{float}{}{\usepackage{float}}
\floatstyle{ruled}
\@ifundefined{c@chapter}{\newfloat{codelisting}{h}{lop}}{\newfloat{codelisting}{h}{lop}[chapter]}
\floatname{codelisting}{Listing}
\newcommand*\listoflistings{\listof{codelisting}{List of Listings}}
\makeatother
\makeatletter
\makeatother
\makeatletter
\@ifpackageloaded{caption}{}{\usepackage{caption}}
\@ifpackageloaded{subcaption}{}{\usepackage{subcaption}}
\makeatother
\usepackage{bookmark}
\IfFileExists{xurl.sty}{\usepackage{xurl}}{} % add URL line breaks if available
\urlstyle{same}
% fallback for those not using the hyperref driver hyperxmp:
\makeatletter
\@ifundefined{xmpquote}{\newcommand{\xmpquote}[1]{#1}}{}
\makeatother
\hypersetup{
  pdftitle={Comments on Orthogonal Knowledge},
  pdfauthor={Brian Weatherson},
  hidelinks,
  pdfcreator={LaTeX via pandoc}}


\title{Comments on \emph{Orthogonal Knowledge}}
\author{Brian Weatherson}
\date{2026}
\begin{document}
\maketitle
\begin{abstract}
Comments at FARM 2026 on Josh Pearson, Richard Roth, and Ginger
Schultheis' \emph{Orthogonal Knowledge}.
\end{abstract}


\setstretch{1.1}
This is a really fun paper, asking a really good question and working
carefully through various possible answers. I'm basically sympathetic to
where they end up at the end of §6. We should give up Inquiry
Orthogonality, and we should resolve the tension in it by accepting
Dimensional Reduction. So I'll just mention one independent reason to
reject Evidential Orthogonality, and then ask a
clarificatory-but-really-critical question about the counterfactual
model in §§7-8.

Evidential Orthogonality says that learning something about different
thermometers can't change what you know about this one. That seems weird
because it doesn't change the probabilities about this thermometer. But
we have independent reason to think that evidence can change what you
know without changing the probability. Here's a version of this I first
learned about from David Chalmers. Say that currently Hero has evidence
\(E_1\) for \(p\), that \(\Pr(p | E_1)\) is really high, but the route
from \(E_1\) to \(p\) goes through some falsehood as in a
Gettier/Dharmottara case, and hence Hero doesn't know \(p\). Then Hero
gets some new evidence that (a) defeats \(E_1\), and substitutes some
new evidence \(E_2\), where \(\Pr(p | E_2) = \Pr(p | E_1)\), but the
route from \(E_2\) to \(p\) is safe. Hero comes to know \(p\) because of
the new evidence, without the probability changing. The thermometer case
isn't the same structure, but it's hard to state a general principle
that entails Evidential Orthogonality without, mistakenly, ruling out
cases like this.

OK, onto the question about the case in §7.2 and §8. What is a world in
these models? Is it an arrangement of material facts, like 000, or does
it settle \emph{all} the facts, including the counterfactuals? To avoid
confusion, call a \textbf{state} a pair like
\(\langle 000, 001\rangle\), where the first member sets the material
facts, and the second says what the material facts are at the nearest
world where they are different. That doesn't settle everything; it
doesn't say what the second nearest material facts are. But it's enough
to raise a problem.

The problem is clearest with different numbers. Say there are 1000
thermometers, each normal with probability 0.99, slightly off with
probability 0.009, and wildly off with probability 0.001. Then an
Ann\(_i\) can know her thermometer is normal, but Carol can't know they
all are. Write \(\langle \mathbf{0}, i\rangle\) for the state where
every thermometer is normal, and the nearest alternative has just
thermometer \(i\) slightly off.

Assume the actual state is \(\langle \mathbf{0}, 1\rangle\), though
Carol doesn't know this. Two questions:

\begin{enumerate}
\def\labelenumi{\arabic{enumi}.}
\tightlist
\item
  What probability should Carol give to
  \(\langle \mathbf{0}, 1\rangle\), and what probability to
  \(\langle \mathbf{0}, 2\rangle\)?
\item
  Where does \(\langle \mathbf{0}, 2\rangle\) come in the ordering of
  states? Is it really close, closer than any state with different
  material facts, as similarity of material facts would suggest? Or is
  it far away, further than states where many thermometers are wildly
  off, as it would be if the ordering is a deep metaphysical fact that's
  harder to change than a thermometer?
\end{enumerate}

I think the answer to 1 has to be that \(\langle \mathbf{0}, 2\rangle\)
is somewhere around 0.001. It's hard to see how Carol could be warranted
in giving it much less probability than
\(\langle \mathbf{0}, 1\rangle\). So let's assume that; if it's wrong,
maybe there's a way out. And with states in the picture, the natural
reading of §7.2 is that Carol knows whatever is true in the first \(m\)
states, where \(m\) is the smallest number making them jointly
\(t\)-likely. So the answer to 2 matters a lot.

Say \(\langle \mathbf{0}, i\rangle\) is really close. The states where
every thermometer is normal are jointly less than \(t\)-likely, so the
first \(m\) states run past them, and every
\(\langle \mathbf{0}, i\rangle\) is among them. By similar reasoning,
the first \(m\) states from \(\langle \mathbf{0}, i\rangle\) include at
least one state where the material facts are different, and at
\(\langle \mathbf{0}, i\rangle\) that means thermometer \(i\) is
slightly off. So at \(\langle \mathbf{0}, i\rangle\), Ann\(_i\) doesn't
know her thermometer is normal, and hence at the actual state no
Ann\(_i\) knows that she knows her thermometer is normal. That holds
however reliable the thermometers are, as long as it's less than
\(t\)-likely that they are all normal.

Say instead \(\langle \mathbf{0}, i\rangle\) is far away. Then the early
part of the ordering carries only a fraction of each arrangement's
probability, roughly a thousandth if the states are about equally
likely. So the first \(m\) states have to reach far down the material
facts, far enough to include states where thermometer \(i\) is not
normal, or even wildly off. Now no Ann\(_i\) even knows that her
thermometer is normal.

Those options aren't exhaustive, and maybe there's a middle way that
avoids both problems. Or maybe it's fine to say that
\(\langle \mathbf{0}, 2\rangle\) has low probability, or that Ann\(_i\)
lacks second-order knowledge. But I'd like to hear more about how the
Authors think probability is distributed over these states, and what
follows from that.

Thanks again to the Authors for a really good paper, asking a great
question and saying some really sharp things about what the answer could
be.



\vspace{12pt} \noindent ~~


\end{document}
